\chapter{Puntos medios y paralelas}\label{ch:g8:midpoints}

Une los puntos medios de dos lados de un triángulo: el \omterm{def:g6:lines:objects}{segmento} que sale es
siempre paralelo al tercer lado y mide exactamente la \omterm{def:g3:division:half}{mitad}. Este
«teorema del punto medio» y su recíproco son el primer aperitivo de una gran
idea --- las \omterm{def:g6:lines:perp}{paralelas} cortan longitudes de forma proporcional --- que florece en el
teorema de Tales del \cref{ch:g9:thales}.

\section{El teorema del punto medio}

\begin{theorem}[Teorema del punto medio]\label{thm:g8:midpoints:direct}
En un triángulo $ABC$, sea $I$ el punto medio de $[AB]$ y $J$ el
punto medio de $[AC]$. Entonces la recta $(IJ)$ es paralela a $(BC)$ y
\[
IJ = \frac{BC}{2} .
\]
\end{theorem}

\begin{proof}[Idea de la demostración]
Sea $K$ el \omterm{def:g7:central:def}{simétrico} de $I$ respecto de $J$ (un medio giro,
\cref{ch:g7:central}). El \omterm{def:g4:geometry:polygon}{cuadrilátero} $AICK$ tiene las diagonales $[IK]$
y $[AC]$ cortándose en su punto medio común $J$: es un
\omterm{def:g7:central:parallelogram}{paralelogramo}, así que $KC$ es paralelo a $AI$, es decir, a $(AB)$, y
$KC = AI = IB$. Entonces $IBCK$ tiene dos lados opuestos ($[IB]$ y
$[KC]$) paralelos e iguales: también es un \omterm{def:g7:central:parallelogram}{paralelogramo}
(\cref{thm:g7:central:recognize}), así que $(IK) \parallel (BC)$ ---
es decir, $(IJ) \parallel (BC)$ --- y $BC = IK = 2\,IJ$.
\end{proof}

\begin{omfigure}
\begin{tikzpicture}[scale=1.1]
  \coordinate (A) at (1.7,2.9);
  \coordinate (B) at (0,0);
  \coordinate (C) at (4.6,0);
  \coordinate (I) at ($(A)!0.5!(B)$);
  \coordinate (J) at ($(A)!0.5!(C)$);
  \draw[thick] (A) -- (B) -- (C) -- cycle;
  \draw[omThm, very thick] (I) -- (J);
  \fill (A) circle (1.5pt) node[above, font=\small] {$A$};
  \fill (B) circle (1.5pt) node[below left, font=\small] {$B$};
  \fill (C) circle (1.5pt) node[below right, font=\small] {$C$};
  \fill[omThm] (I) circle (1.5pt) node[left, font=\small] {$I$};
  \fill[omThm] (J) circle (1.5pt) node[right, font=\small] {$J$};
  \draw ($(A)!0.22!(B)$) ++(-0.09,-0.05) -- ++(0.18,0.1);
  \draw ($(A)!0.72!(B)$) ++(-0.09,-0.05) -- ++(0.18,0.1);
  \draw ($(A)!0.25!(C)$) ++(-0.06,-0.1) -- ++(0.12,0.2);
  \draw ($(A)!0.25!(C)$) ++(0.02,-0.1) -- ++(0.12,0.2);
  \draw ($(A)!0.75!(C)$) ++(-0.06,-0.1) -- ++(0.12,0.2);
  \draw ($(A)!0.75!(C)$) ++(0.02,-0.1) -- ++(0.12,0.2);
\end{tikzpicture}

{\small El \omterm{def:g6:lines:objects}{segmento} de los puntos medios $[IJ]$ (en rojo) es paralelo al tercer
lado $(BC)$ y mide la \omterm{def:g3:division:half}{mitad}.}
\end{omfigure}

\begin{example}\label{ex:g8:midpoints:direct}
En un triángulo $ABC$ con $BC = 9$ cm, se unen los puntos medios $I$ de $[AB]$ y
$J$ de $[AC]$. Sin medir nada:
$(IJ) \parallel (BC)$ y $IJ = \frac92 = 4.5$ cm. El triángulo pequeño
$AIJ$ es una copia a \omterm{def:g3:division:half}{mitad} de tamaño de $ABC$ --- su \omterm{def:g6:measure:perimeter}{perímetro} es también la \omterm{def:g3:division:half}{mitad}
del \omterm{def:g6:measure:perimeter}{perímetro} de $ABC$.
\end{example}

\begin{theorem}[Recíproco]\label{thm:g8:midpoints:converse}
En un triángulo $ABC$, la recta que pasa por el punto medio $I$ de $[AB]$ y es
\emph{paralela} a $(BC)$ corta al lado $[AC]$ en su punto medio.
\end{theorem}
\admitted

\begin{example}\label{ex:g8:midpoints:converse}
Un camino atraviesa un campo triangular partiendo de la \omterm{def:g3:division:half}{mitad} de un borde
y siguiendo paralelo al borde opuesto. El recíproco garantiza
que el camino sale exactamente por la \omterm{def:g3:division:half}{mitad} del otro borde --- sin
necesidad de medir al otro lado.
\end{example}

\begin{method}[Elegir entre el teorema y el recíproco]\label{met:g8:midpoints:choose}
\begin{enumerate}
  \item \emph{dos puntos medios conocidos} $\Rightarrow$ usa el teorema
        directo: concluye paralelismo y \omterm{def:g3:division:half}{mitad} de longitud;
  \item \emph{un punto medio y una paralela conocidos} $\Rightarrow$ usa
        el recíproco: concluye que el punto de corte es un punto medio;
  \item escribe qué triángulo, qué puntos medios y qué teorema
        usas --- una frase para cada cosa.
\end{enumerate}
\end{method}

\section{Triángulos a mitad de tamaño}

\begin{proposition}[El triángulo de los puntos medios]\label{prop:g8:midpoints:medial}
Unir los tres puntos medios de los lados de un triángulo lo corta en
cuatro triángulos de \omterm{def:g6:measure:area}{áreas} iguales, cada uno copia a \omterm{def:g3:division:half}{mitad} de tamaño del original.
\end{proposition}

\begin{proof}
Cada lado del triángulo de los puntos medios es paralelo a un lado de $ABC$ y mide la
\omterm{def:g3:division:half}{mitad} (\cref{thm:g8:midpoints:direct}, aplicado tres veces).
Los cuatro triángulos pequeños tienen lados de las mismas tres longitudes, así que
son copias idénticas; juntos recubren $ABC$, así que cada uno tiene un
\omterm{def:g3:division:half}{cuarto} de su \omterm{def:g6:measure:area}{área}.
\end{proof}

\begin{omfigure}
\begin{tikzpicture}[scale=1.05]
  \coordinate (A) at (1.9,3.0);
  \coordinate (B) at (0,0);
  \coordinate (C) at (5.0,0);
  \coordinate (I) at ($(A)!0.5!(B)$);
  \coordinate (J) at ($(A)!0.5!(C)$);
  \coordinate (K) at ($(B)!0.5!(C)$);
  \draw[thick] (A) -- (B) -- (C) -- cycle;
  \fill[omThm!12] (I) -- (J) -- (K) -- cycle;
  \draw[omThm, thick] (I) -- (J) -- (K) -- cycle;
  \fill (A) circle (1.4pt) node[above, font=\small] {$A$};
  \fill (B) circle (1.4pt) node[below left, font=\small] {$B$};
  \fill (C) circle (1.4pt) node[below right, font=\small] {$C$};
  \fill[omThm] (I) circle (1.4pt) node[left, font=\small] {$I$};
  \fill[omThm] (J) circle (1.4pt) node[right, font=\small] {$J$};
  \fill[omThm] (K) circle (1.4pt) node[below, font=\small] {$K$};
\end{tikzpicture}

{\small El triángulo de los puntos medios $IJK$ corta $ABC$ en cuatro triángulos
iguales --- una imagen que conviene recordar.}
\end{omfigure}

\begin{remark}[Camino de Tales]
El teorema del punto medio es el caso «una \omterm{def:g3:division:half}{mitad}» de un hecho más general:
una recta paralela a un lado de un triángulo corta los otros dos lados en
\emph{razones iguales} --- un tercio y un tercio, dos quintos y dos
quintos\,\dots\ Ese enunciado general es el teorema de Tales
(\cref{ch:g9:thales}); el caso del punto medio es el único que hace falta este
año.
\end{remark}

\section{Ejercicios}

\begin{exercise}[$\star$]\label{exo:g8:midpoints:1}
En un triángulo $RST$, $M$ es el punto medio de $[RS]$ y $N$ el
punto medio de $[RT]$, con $ST = 7$ cm. ¿Qué se puede decir de la recta
$(MN)$ y de la longitud $MN$? Cita el teorema que usas.
\end{exercise}

\begin{exercise}[$\star$]\label{exo:g8:midpoints:2}
En un triángulo $ABC$, $I$ es el punto medio de $[AB]$ y la recta
que pasa por $I$ paralela a $(BC)$ corta a $[AC]$ en $J$, con
$AC = 11$ cm. Calcula $AJ$ citando el teorema que usas.
\end{exercise}

\begin{exercise}[$\star$]\label{exo:g8:midpoints:3}
Dibuja un triángulo $ABC$ con $AB = 8$ cm, $AC = 6$ cm y $BC = 7$ cm,
sitúa los puntos medios $I$ de $[AB]$ y $J$ de $[AC]$, y mide $IJ$
para comprobar el teorema. Calcula el \omterm{def:g6:measure:perimeter}{perímetro} de $AIJ$ sin
medir nada más.
\end{exercise}

\begin{exercise}[$\star$]\label{exo:g8:midpoints:4}
$IJK$ es el triángulo de los puntos medios de $ABC$, cuyos lados miden $10$,
$12$ y $16$ cm. Da los tres lados de $IJK$ y su \omterm{def:g6:measure:perimeter}{perímetro}.
¿Cómo se comparan los dos \omterm{def:g6:measure:perimeter}{perímetros}?
\end{exercise}

\begin{exercise}[$\star$]\label{exo:g8:midpoints:5}
El \omterm{def:g6:measure:area}{área} de un triángulo es $36$ cm$^2$. ¿Cuál es el \omterm{def:g6:measure:area}{área} de su
triángulo de los puntos medios? ¿Y la de cada uno de los cuatro triángulos pequeños?
\end{exercise}

\begin{exercise}[$\star\star$]\label{exo:g8:midpoints:6}
En un triángulo $ABC$, $I$ es el punto medio de $[AB]$, $J$ el de
$[AC]$ y $K$ el de $[BC]$. Demuestra que $IJKB$ es un \omterm{def:g7:central:parallelogram}{paralelogramo}.
(Compara $[IJ]$ y $[BK]$: ¿paralelos? ¿iguales?)
\end{exercise}

\begin{exercise}[$\star\star$]\label{exo:g8:midpoints:7}
Una vela triangular tiene el borde inferior de $6.4$ m. Una costura de refuerzo
une los puntos medios de los otros dos bordes. ¿Qué longitud de costura
hace falta? ¿Y si la costura uniera en cambio los puntos situados a un
\omterm{def:g3:division:half}{cuarto} de cada borde desde el \omterm{def:g5:solids:def}{vértice} superior? (Conjetura con
un dibujo; la demostración es de Tales, \cref{ch:g9:thales}.)
\end{exercise}

\begin{exercise}[$\star\star$]\label{exo:g8:midpoints:8}
En el triángulo $ABC$, rectángulo en $A$, $M$ es el punto medio de la
hipotenusa $[BC]$, y la paralela a $(AB)$ por $M$ corta a
$[AC]$ en $N$.
\begin{enumerate}
  \item demuestra que $N$ es el punto medio de $[AC]$;
  \item explica por qué $(MN)$ es perpendicular a $(AC)$.
\end{enumerate}
\end{exercise}

\begin{exercise}[$\star\star\star$]\label{exo:g8:midpoints:9}
Sea $ABCD$ un \omterm{def:g4:geometry:polygon}{cuadrilátero} \emph{cualquiera} y sean $I$, $J$, $K$ y $L$ los
puntos medios de $[AB]$, $[BC]$, $[CD]$ y $[DA]$ (la conjetura del
\cref{exo:g7:central:11}).
\begin{enumerate}
  \item aplica el teorema del punto medio en el triángulo $ABC$ al
        \omterm{def:g6:lines:objects}{segmento} $[IJ]$, y en el triángulo $ACD$ al \omterm{def:g6:lines:objects}{segmento}
        $[LK]$: compara cada uno con la diagonal $[AC]$;
  \item concluye que $IJKL$ es siempre un \omterm{def:g7:central:parallelogram}{paralelogramo}.
\end{enumerate}
\end{exercise}

\section{Problema: Las tres medianas y el centro de gravedad}

\begin{problem}[{Problema de fin de semana --- las medianas de un triángulo se cortan
en un solo punto, a dos tercios del recorrido de cada una}]
\label{pb:g8:midpoints:1}
Recorta un triángulo de cartón rígido y se sostendrá en equilibrio,
perfectamente plano, sobre la punta de un lápiz colocada en un punto especial: el
\emph{centro de gravedad}\index{centro de gravedad} del triángulo.
Este problema halla ese punto. Una \emph{mediana}\index{mediana} de un
triángulo es el \omterm{def:g6:lines:objects}{segmento} que une un \omterm{def:g5:solids:def}{vértice} con el \emph{punto medio} del
lado opuesto. Demostrarás --- con el teorema del punto medio
usado dos veces de una manera inesperada --- que las tres medianas pasan
todas por un mismo punto, y lo localizarás con precisión.

\medskip
\noindent\textbf{Parte I --- Calentamiento.}
\begin{enumerate}
  \item dibuja un triángulo grande $ABC$ (sin forma especial), construye
        los puntos medios $I$ de $[AB]$, $J$ de $[AC]$ y $K$ de
        $[BC]$, y traza las tres medianas $[AK]$, $[BJ]$ y $[CI]$.
        ¿Qué observas?
  \item demuestra que una mediana corta el triángulo en dos triángulos
        de \omterm{def:g6:measure:area}{áreas} iguales. (Compara bases y alturas, y usa la
        fórmula del \omterm{def:g6:measure:area}{área}: base $\times$ altura $\div\ 2$.) Para un
        triángulo de \omterm{def:g6:measure:area}{área} $36$ cm$^2$, ¿cuáles son las \omterm{def:g6:measure:area}{áreas} de los
        dos trozos?
  \item en el triángulo $ABC$, ¿qué dice
        el \cref{thm:g8:midpoints:direct} sobre el \omterm{def:g6:lines:objects}{segmento}
        $[IJ]$?
\end{enumerate}

\medskip
\noindent\textbf{Parte II --- Dos medianas se cortan a dos tercios del
recorrido.}
Las medianas $[BJ]$ y $[CI]$ se cortan en un punto; llámalo $G$. Sea
$M$ el punto medio de $[GB]$ y $N$ el punto medio de $[GC]$.
\begin{enumerate}[resume]
  \item aplica el teorema del punto medio \emph{en el triángulo $GBC$}:
        ¿qué dice sobre el \omterm{def:g6:lines:objects}{segmento} $[MN]$?
  \item deduce que $(IJ) \parallel (MN)$ y $IJ = MN$, y
        concluye con el \cref{thm:g7:central:recognize} que $IJNM$
        es un \omterm{def:g7:central:parallelogram}{paralelogramo};
  \item explica por qué $I$ y $N$ están los dos sobre la mediana $(CI)$,
        y $J$ y $M$ los dos sobre la mediana $(BJ)$ --- así que las
        diagonales del \omterm{def:g7:central:parallelogram}{paralelogramo} $IJNM$ son $[IN]$ y
        $[JM]$, y se cortan exactamente en $G$. ¿Qué dice entonces la
        \emph{definición} de \omterm{def:g7:central:parallelogram}{paralelogramo} sobre $G$?
  \item deduce que $BM = MG = GJ$ y concluye:
        \[
        BG = 2 \times GJ
        \]
        --- el punto $G$ está sobre la mediana $[BJ]$ a dos tercios
        del recorrido desde el \omterm{def:g5:solids:def}{vértice} $B$. Enuncia y justifica el
        resultado análogo sobre la mediana $[CI]$;
  \item olvídate ahora de $[CI]$ y repite exactamente el mismo argumento con la
        pareja de medianas $[BJ]$ y $[AK]$: su punto de corte
        también está a dos tercios de $[BJ]$ desde $B$. Explica por qué
        eso obliga a que sea el \emph{mismo} punto $G$ --- y
        concluye que las tres medianas pasan por $G$;
  \item en un triángulo cuya mediana $[BJ]$ mide $9$ cm,
        ¿a qué distancia está $G$ de $B$ y de $J$?
  \item justifica (en una frase) que también $AG = 2 \times GK$.
\end{enumerate}

\medskip
\noindent\textbf{Parte III --- Seis porciones iguales y coordenadas.}
Las tres medianas cortan el triángulo en seis triángulos pequeños
alrededor de $G$.
\begin{enumerate}[resume]
  \item en el triángulo $GBC$, el \omterm{def:g6:lines:objects}{segmento} $[GK]$ es una mediana.
        Deduce que los dos triángulos pequeños $GBK$ y $GKC$ tienen
        \omterm{def:g6:measure:area}{áreas} iguales;
  \item usando la pregunta 2 en los triángulos $ABK$ y $AKC$ (los dos
        cortados por la mediana $[AK]$ de $ABC$), demuestra que los
        triángulos $ABG$ y $ACG$ tienen \omterm{def:g6:measure:area}{áreas} iguales. Repitiendo
        con otra mediana, concluye que los tres triángulos
        $ABG$, $BCG$ y $CAG$ tienen todos \omterm{def:g6:measure:area}{área} un tercio de la de $ABC$;
  \item cada uno de esos tres triángulos queda cortado por la \omterm{def:g3:division:half}{mitad} por un trozo
        de mediana (por ejemplo, $[GI]$ es una mediana del
        triángulo $ABG$ --- ¿vista desde qué \omterm{def:g5:solids:def}{vértice}?). Concluye: las seis
        porciones pequeñas de alrededor de $G$ tienen todas la misma \omterm{def:g6:measure:area}{área}, una sexta parte
        del triángulo.
  \item en una cuadrícula de coordenadas, sitúa $A(0, 0)$, $B(6, 0)$ y
        $C(0, 6)$, con $K(3, 3)$ el punto medio de $[BC]$. Usa el
        teorema de los dos tercios de la Parte II para calcular las coordenadas
        de $G$ sobre la mediana $[AK]$; comprueba después que el mismo
        punto está a dos tercios de la mediana $[BJ]$, donde
        $J(0, 3)$. Compara las coordenadas de $G$ con las medias
        \[
        \frac{0 + 6 + 0}{3},
        \qquad
        \frac{0 + 0 + 6}{3} .
        \]
  \item el final del equilibrio: usando las preguntas 2 y 8, explica por qué
        el triángulo de cartón se sostiene sobre el filo de un cuchillo colocado
        a lo largo de cualquier mediana --- la misma cantidad de cartón a cada
        lado --- y por qué, en consecuencia, la punta del lápiz tiene que colocarse
        en $G$, el punto que los físicos llaman \omterm{pb:g8:midpoints:1}{centro de
        gravedad}. (Una demostración del equilibrio totalmente rigurosa necesita el
        cálculo integral de los volúmenes universitarios; las \omterm{def:g6:measure:area}{áreas}
        iguales lo hacen creíble hoy.)
\end{enumerate}
\end{problem}
